\textbf{Hipótese proposta:}
O pombo-torcaz e o urubu possuem estratégias de voo e compromissos morfológicos fundamentalmente diferentes, o que resulta em frações de massa muscular $W_{fm}/W_b$ distintas.

\textbf{Argumento de escala:}
\begin{itemize}
\item \textbf{Pombo-torcaz (\textit{Columba palumbus}):} batedor ativo, com frequência de batenção de asas elevada. Para gerar sustentação e propulsão mecânica por batimento, precisa de muscúlatura de voo com grande área de secção transversal (força muscular $\propto L^2$). Isso resulta em \textbf{alta razão $W_{fm}/W_b$}.
\item \textbf{Urubu (\textit{Cathartes aura}):} planador térmico. Usa correntes ascendentes de ar quente para manter altitude sem bater as asas. A carga alar baixa ($W/S$, peso por unidade de área alar) e a eficiência aerodinâmica do planejo reduzem a demanda sobre os músculos de voo, resultando em \textbf{baixa razão $W_{fm}/W_b$}.
\end{itemize}

\textbf{Conclusão:} as diferenças na Figura 3.2 refletem nicho ecológico e modo de voo, não apenas tamanho corporal. A escala simples $W_{fm} \propto W_b$ só é válida para espécies com estratégias de voo homogêneas.

